% Generated from paper/full_draft. Edit the Markdown source or migration script.

\section{Discussion}
\label{sec:07-discussion:discussion}

\subsection{Supported Claim}
\label{sec:07-discussion:supported-claim}

IntentWeight supports a bounded but useful claim. It is a feedback-driven adaptive evidence-selection controller instantiated as a retrieval controller in the evaluated QA setting. It can control final context budget while preserving dense-level retrieval quality. On LoTTE technology/search, the conservative confidence-based context policy reduces final retrieved context tokens by about 4.7-5.3\% from 100k to 638k corpus chunks. Mean $\mathrm{Hit@10}$ is above dense-only retrieval on 200k, 400k, and 638k.

This result is not a claim that dense retrieval is weak. Dense retrieval remains the primary quality baseline and an important recall floor. The value of IntentWeight is that it learns when dense fallback is needed, when local geometry is reliable, and when the final context can be safely compacted.

\subsection{Why Multi-Route Retrieval Alone Is Not Enough}
\label{sec:07-discussion:why-multi-route-retrieval-alone-is-not-enough}

A static combination of dense, BM25, and cluster-local retrieval can improve coverage, but it does not automatically reduce final context tokens. In fact, static dense+BM25 hybrid retrieval can use more context tokens than dense-only retrieval because it surfaces longer or noisier chunks. The token-saving mechanism is therefore not "more routes." It is confidence-based final context control.

The main confidence-based policy is intentionally conservative. It compresses only high-confidence cases to $k=8$ and keeps mid-confidence cases at $k=10$. This is why the saving is modest but stable. More aggressive configurations save more tokens, but they produce visible $\mathrm{Hit@10}$ loss. The conservative policy should therefore be interpreted as a safe operating point on a token-quality frontier.

\subsection{Feedback Improves the Policy Field}
\label{sec:07-discussion:feedback-improves-the-policy-field}

Dense and BM25 fallback can saturate final $\mathrm{Hit@10}$. This can hide the effect of feedback in final fused retrieval metrics. The clearer feedback signal appears in route-policy metrics such as selected-cluster hit and last true reward.

Trust weighting improves these policy metrics under controlled simulated feedback. This supports the idea that the system can self-improve under a usable feedback signal. However, the current feedback is simulated and ground-truth-derived. Production systems still need real feedback collection, trust scoring, delayed-feedback handling, and safeguards against unreliable or adversarial signals.

\subsection{Geometry Is Useful but Not Sufficient}
\label{sec:07-discussion:geometry-is-useful-but-not-sufficient}

The geometry diagnostics support a piecewise relevance-manifold framing. $\mathrm{NearestClusterHit@3}$ remains high across LoTTE scales, and local geometry provides useful routing information. However, context retention declines with scale, and geometry alone is not a complete retrieval model. If a cluster route prunes too early, correct evidence can be lost.

IntentWeight therefore uses geometry as one signal in a controller. Dense retrieval remains a fallback, BM25 provides lexical anchors, and LinUCB learns route confidence over repeated interactions.

\subsection{Evidence Completeness Versus Usable Evidence}
\label{sec:07-discussion:evidence-completeness-versus-usable-evidence}

The main retrieval headline is query-level $\mathrm{Hit@10}$. This metric asks whether at least one relevant chunk appears in the final context. It is appropriate for many RAG settings where one good supporting chunk is enough to ground an answer. However, context compaction can reduce $\mathrm{EvidenceRecall@10}$, the fraction of all GT chunks retrieved.

This is an expected trade-off. IntentWeight optimizes usable evidence under a smaller context budget, not exhaustive evidence collection. For legal review, medical synthesis, or compliance workflows where complete evidence coverage is required, the system should use a more conservative context policy or disable compaction.

\subsection{Production Interpretation}
\label{sec:07-discussion:production-interpretation}

The measured 4.7-5.3\% token reduction is conservative. It is measured with limited simulated interaction history and a cautious $k=8$ high-confidence compression policy. In production, repeated query patterns, user-specific feedback, and richer confidence tiers may increase the fraction of queries that can be safely compacted. This is a hypothesis for production evaluation rather than a claim proven by the current experiments.

The correct deployment interpretation is therefore:

\begin{itemize}
\item keep dense retrieval as a recall floor;
\item use feedback and confidence to reduce final context when the policy is stable;
\item monitor evidence quality and fallback rates;
\item avoid aggressive compaction for complete-evidence tasks;
\item treat token saving as a controllable frontier rather than a fixed guarantee.
\end{itemize}
